\section{Validation Methods}\label{sec:methods}

The analysis plan was written on 29~August~2026, when zero races were in
hand, and is kept in the repository with an amendment log. Everything in
this section was fixed before any parameter was fitted to any race. The
plan's function is to make the empirical section falsifiable in advance:
what is measured, what counts as a win, what is fitted, and what is done
about the dive start were all decided when the data could not yet influence
the decision.

\subsection{Primary metric}

For each race with cumulative splits $S_1<S_2<S_3<S_4=T$ we form lap times
$s_i=S_i-S_{i-1}$ and the free-swimming-equivalent shares
\begin{equation}
  P_i \;=\; \frac{s_i + \delta_{i1}\,c}{T + c}, \qquad i=1,\dots,4,
  \label{eq:shares}
\end{equation}
where $c$ is the dive-start credit added back to lap~1
(Section~\ref{sec:start}). The model predicts raced-space fractions
$P_i^\ast$, and the per-race primary metric is
$\mathrm{RMSE}=\big[\tfrac14\sum_i (P_i-P_i^\ast)^2\big]^{1/2}$, reported
in percentage points and averaged over races. The metric is scale-free, so a
1:37 swimmer and a 2:30 swimmer contribute comparably. Because the $P_i$ sum
to one, only three residuals are free; RMSE over four terms remains a valid
comparison statistic since every model is scored identically, but it is
never treated as having four independent components. Secondary metrics are
the mean absolute error, the per-lap signed error (which is what actually
separates M3 from M4), and the pre-registered discriminator: the lap
$1\!\to\!2$ drop against the lap $3\!\to\!4$ drop. Race-time error is
reported but never used for selection, because race time depends on
parameters the shape cannot identify (Theorem~\ref{thm:shape}). Likelihood
criteria such as AIC are not used: the models are deterministic and carry no
error model, and inventing one for the convenience of a statistic would be a
modelling choice made backwards.

\subsection{What counts as a win}

A model \emph{wins on shape} if its predicted sign of $P_4-P_1$ matches the
observed mean's and no other model's does; M0 and M1 are shape-identical and
count once. A model \emph{wins on accuracy} if its held-out mean RMSE is
lower than every other model's by more than the 95\% bootstrap confidence
half-width of that difference. If neither criterion separates the models,
the registered conclusion is that the data do not distinguish them, and
that is reported as the result.

\subsection{Split, fitting, and the leakage guard}

Races are split into training and test sets \emph{by swimmer}, never by race:
repeated races from one swimmer share technique, start quality and habitual
pacing, so a race-level split would score models partly on how well they
memorized individuals. Swimmers are assigned whole to one side until the
test side holds at least a quarter of the races, with a seed (20260829)
fixed in the plan. On pilot-v0.2 this gives 256 training races by 188
swimmers and 89 held-out races by 60 swimmers. Fitting code can obtain data
only through a function that withholds the test rows, and a structural guard
raises an exception if any frame containing a held-out swimmer reaches an
objective; the guard is unit-tested. The test set is opened once, by one
script, which records the time it did so.

Exactly one parameter per model is fitted to shape: $\beta_E$ for M2,
$\beta_x$ for M3, $\gamma$ for M4; M0 and M1 have none. Nothing else is
fitted to pacing data. By Theorem~\ref{thm:shape} the engine parameters
($E_0$, $R$, $\tau$), the drag block and the efficiencies do not enter the
optimal shape, so fitting them to splits would produce numbers that are
pure noise, and reporting those as physiology would be the worst error
available to this project. Each fit minimizes the mean training RMSE over
its parameter: for M3 the shape is closed-form and a bounded scalar search
is exact; for M2 and M4 every objective evaluation is an inner
optimal-control solve, so a coarse grid brackets the minimum and a bounded
Brent refinement follows, with the evaluation budget recorded. One free
parameter each makes the comparison fair without a complexity penalty.

\subsection{The dive start}\label{sec:start}

The model describes free swimming; recorded lap~1 contains a dive. Because
the start and both positive-split couplings push lap~1 the same way, an
uncorrected fit would absorb the dive into the fatigue parameter and
overstate it. The plan therefore (i) adds a credit $c$ to lap~1 before
computing shares, with $c=1.80$~s anchored to measured elite 15~m start
times \citep{tor2014characteristics,rudnik2023kinematic}; (ii) reports every
headline result across a sensitivity band, originally 1.2--2.4~s and
widened by two logged amendments to 1.2--3.4~s once those same start times
were converted to this field's slower race pace
($c \approx 15/\bar v - t_{15} \approx 3.1$--$3.4$~s); and (iii) rules that
if the model ranking changes anywhere inside the band, the conclusion is
that the data cannot distinguish the models given start uncertainty. A
third amendment, logged on 1~September~2026, corrected a sign error in the
implementation of item~(i): the credit had been subtracted from observed
lap~1 as well as from the model's, double-counting it. The metric definition
did not change; the code was made to match it, and every pilot output was
regenerated. All three amendments preceded any calibration and any use of
the test set, and appear in the plan's amendment log rather than as silent
edits.

\subsection{Deviation and performance}

The primary hypothesis is that swimmers whose pacing lies closer to the
predicted optimum perform better relative to expectation. For race $j$ of
swimmer $i$ we compute the deviation $D_{ij}$, the RMSE of the race's shares
from the best-supported model's fitted optimum, and the relative improvement
$I_{ij}=(T^{\mathrm{PB}}_{ij}-T_{ij})/T^{\mathrm{PB}}_{ij}$, where
$T^{\mathrm{PB}}_{ij}$ is the swimmer's best time \emph{before} that race; a
later best would leak the outcome into the predictor. The registered model is
a mixed-effects regression with a random intercept per swimmer,
$I_{ij}=b_0+b_1 D_{ij}+b_2 D_{ij}^2+u_i+e_{ij}$, quadratic because the
theory predicts a quadratic penalty near the optimum
(Section~\ref{sec:results-theory}), with a pre-registered fallback to
ordinary least squares with swimmer-clustered standard errors if the
random-effect variance is not identifiable at the achieved sample size.
Confidence intervals for every held-out quantity are percentile bootstrap
intervals that resample swimmers, not races (10{,}000 draws, one resample
applied to all models so that differences are computed on matched draws).

\subsection{Power, stated in advance}

The theoretical effect is small: a 5\% opening error costs about
0.09~s in 100~s, while meet-to-meet variation in a swimmer's 200 free is of
order 1\%. The plan states that even 200--300 races may be underpowered for
the deviation--performance relationship and that a null there will be
reported as a null. The model-comparison question is better powered, because
the shape differences between M2, M3 and M4 are large relative to the
race-to-race noise in shares.

\subsection{Exploratory versus registered}

The 80-race first pilot was used to exercise the fitting machinery and to
quantify the start-credit confound; those fits are labelled exploratory in
the repository and are not the calibration. The registered calibration was
run on the training side of the frozen pilot-v0.2 dataset after the freeze
(with one inner-solver fix, described in
Section~\ref{sec:results-empirical}, made and logged before the test set
was opened); the held-out evaluation was then run once on its test side. All code, the plan, its amendments, the dataset
register with source digests, and a continuous-integration test suite are in
the project repository. The anonymized race table itself is not published,
because ages, teams and exact times at named meets would re-identify
minors; the chain from official result sheet to every number in this paper
can nonetheless be re-executed by anyone holding those public sheets, whose
digests are recorded.
